% !TEX root = ../../main.tex
% Layout inspiration: Dave Richeson's LaTeX cheat sheet,
% divisbyzero.com, Dickinson College, as credited in the supplied example.
% Adapted to the book's portrait pages with two columns and analysis formulas.
\begingroup
\titlespacing*{\chapter}{0pt}{0pt}{14pt}
\chapter{Formula sheet}
\label{ch:appendix:formula-sheet}

\small
\setlength{\parindent}{0pt}
\setlength{\columnsep}{24pt}
\setlength{\multicolsep}{8pt}
\setlength{\abovedisplayskip}{5pt plus 2pt minus 1pt}
\setlength{\belowdisplayskip}{5pt plus 2pt minus 1pt}
\setlength{\abovedisplayshortskip}{3pt plus 1pt}
\setlength{\belowdisplayshortskip}{5pt plus 2pt minus 1pt}
\titlespacing*{\section}{0pt}{0pt}{9pt}
\newcommand{\formulagroup}[1]{%
  \par\addvspace{9pt}\noindent\textbf{#1}\par\nobreak\smallskip
}
% Side limits keep the reference sums compact without reducing the font.
\newcommand{\sheetsum}{\sum\nolimits}

\begin{fullwidth}
  Throughout, limits are taken as the index tends to infinity and $N$ is
  a positive integer.

  \begin{multicols}{2}
    \raggedcolumns
    \section{Sequences}
    \label{sec:appendix:formula-sheet-sequences}

    \formulagroup{Convergence}
    The condition $x_n\to L$ means that for every $\varepsilon>0$ there
    exists $N\in\N$ such that
    \[
      \forall n\geq N\quad |x_n-L|<\varepsilon.
    \]

    \formulagroup{Limit laws}
    If $x_n\to L$, $y_n\to M$, and $a,b\in\R$, then
    \[
      \begin{aligned}
        ax_n+by_n&\longrightarrow aL+bM,\\
        x_ny_n&\longrightarrow LM,\\
        \frac{x_n}{y_n}&\longrightarrow\frac{L}{M}
          \quad(M\ne0).
      \end{aligned}
    \]

    \formulagroup{Squeeze theorem}
    If $u_n\leq x_n\leq v_n$ eventually, then
    \[
      u_n\to L,\quad v_n\to L
      \quad\Longrightarrow\quad x_n\to L.
    \]

    \formulagroup{Bounded monotone sequences}
    If $(x_n)$ is bounded and monotone, then
    \[
      \lim x_n=
      \begin{cases}
        \sup_{n\geq1}x_n & \text{if nondecreasing},\\
        \inf_{n\geq1}x_n & \text{if nonincreasing}.
      \end{cases}
    \]

    \formulagroup{Cauchy criterion}
    A real sequence converges if and only if for every $\varepsilon>0$
    there exists $N\in\N$ such that
    \[
      \forall m,n\geq N\quad |x_m-x_n|<\varepsilon.
    \]

    \formulagroup{Subsequences and limit points}
    A subsequence has indices $n_1<n_2<\cdots$.
    If $x_n\to L$, then $x_{n_k}\to L$.
    A limit point is the limit of some subsequence.

    \formulagroup{Upper and lower limits}
    For a bounded real sequence,
    \[
      \begin{aligned}
        \liminf x_n&=\lim_{n\to\infty}\inf\{x_k\mid k\geq n\},\\
        \limsup x_n&=\lim_{n\to\infty}\sup\{x_k\mid k\geq n\}.
      \end{aligned}
    \]
    These are the least and greatest subsequence limits.
    \[
      x_n\to L\quad\Longleftrightarrow\quad
      \liminf x_n=\limsup x_n=L.
    \]

    \columnbreak
    \section{Series}
    \label{sec:appendix:formula-sheet-series}

    \formulagroup{Partial sums and convergence}
    \[
      S_N=\sheetsum_{n=1}^{N}a_n,\qquad
      \sheetsum_{n=1}^{\infty}a_n=S\ \Longleftrightarrow\ S_N\to S.
    \]
    Convergence requires $a_n\to0$. The converse fails.

    \formulagroup{Finite geometric sum}
    For $r\ne1$,
    \[
      \sheetsum_{k=0}^{N}ar^k=a\frac{1-r^{N+1}}{1-r}.
    \]
    When $r=1$, the sum is $a(N+1)$.

    \formulagroup{Geometric series and remainder}
    If $|r|<1$, then
    \[
      \begin{aligned}
        \sheetsum_{k=0}^{\infty}ar^k&=\frac{a}{1-r},\\
        \sheetsum_{k=N+1}^{\infty}ar^k&=\frac{ar^{N+1}}{1-r}.
      \end{aligned}
    \]

    \formulagroup{Telescoping sums}
    \[
      \sheetsum_{n=1}^{N}(b_n-b_{n+1})=b_1-b_{N+1}.
    \]
    If $b_n\to B$, it follows that
    \[
      \sheetsum_{n=1}^{\infty}(b_n-b_{n+1})=b_1-B.
    \]

    \formulagroup{Arithmetic sums}
    For a progression with first term $a$ and difference $d$,
    \[
      \sheetsum_{k=1}^{N}\bigl(a+(k-1)d\bigr)
      =\frac{N}{2}\bigl(2a+(N-1)d\bigr).
    \]
    \[
      \begin{aligned}
        \sheetsum_{k=1}^{N}k&=\frac{N(N+1)}{2},\\
        \sheetsum_{k=1}^{N}k^2&=\frac{N(N+1)(2N+1)}{6},\\
        \sheetsum_{k=1}^{N}k^3&=\left(\frac{N(N+1)}{2}\right)^2.
      \end{aligned}
    \]

    \formulagroup{Two familiar series}
    \[
      \sheetsum_{n=0}^{\infty}\frac{1}{n!}=e,\qquad
      \sheetsum_{n=1}^{\infty}\frac{1}{n(n+1)}=1.
    \]
  \end{multicols}
\end{fullwidth}
\endgroup
